\section{Invariant bases}
\label{app:bases}

\paragraph{Identifiers.}
Candidates carry stable identifiers combining a serial number, a family name and
a field degree, as in \texttt{c046\_portgraph\_d12}. The family is one of
\texttt{portgraph}, \texttt{structured} or \texttt{trace}. Identifiers are fixed once
and never reused, so a candidate referenced in a certificate can always be
located. The candidate ordering used in the rank computation is itself hashed and
the hash is recorded with every cell, which is what allows cells produced in
different working trees to be compared at all.

\paragraph{Degree-four and degree-six.}
$A_4$ is one-dimensional, spanned by $\operatorname{Tr}(M^2)$ in the normalisation
of appendix~\ref{app:blocks}, equivalently by $I^{(4)}_1$. $A_6$ is
two-dimensional; a convenient basis is $\{\operatorname{Tr}(M^3), I^{(6)}_1\}$.

\paragraph{Degree eight.}
$\dim A_8 = \dEightUnionRank$. The tensor-side basis consists of trace monomials
in $M$ together with contractions involving the quartic blocks. On the spinor side
the corresponding family requires the tensor words discussed in
section~\ref{sec:deg8}: without them the family spans only
$\dEightRankWithoutWords$ dimensions.

\paragraph{Degree ten.}
$\dim_\Q \Ften = \dimAtenQ$, decomposing as
\begin{equation}
\Ften = \Gten \oplus \Pten,
\qquad
\dim \Gten = \dimGten,
\qquad
\dim \Pten = \dimPten .
\end{equation}
Three subspaces of $\Ften$ recur:
\begin{itemize}\itemsep2pt
\item $\Bten$, the span of the published structures, $\dim_\Q \Bten = \dimBten$;
\item $\Dten$, the flow-reachable subspace, $\dim_\Q \Dten = \dimDtenQ$;
\item $\Pten$, the products, $\dim_\Q \Pten = \dimPten$.
\end{itemize}
Their pairwise incidences are tabulated in table~\ref{tab:spans}.

\paragraph{The $\Omega$ and $\Theta$ candidates.}
Two families on the spinor side are written $\Omega$ and $\Theta$; they are the
tensor words of section~\ref{sec:spinor}, distinguished by whether the
intermediate tensor structure is quadratic or quartic. They are listed
separately because the degree-eight ablation removes exactly these.

\paragraph{The ambiguous published candidates.}
Two of the published degree-ten expressions cannot be read unambiguously from the
source PDF, because bracket colour --- which fixes the order of operations ---
does not survive text extraction. Both readings of each are implemented, denoted
by an \texttt{-a} or \texttt{-b} suffix on the identifier, and both are projected
and reported. The compact $\Qten$ basis is deliberately chosen to avoid these
candidates, so no result in this paper depends on which reading is correct. This
is an avoidance, not a resolution; see section~\ref{sec:limitations}.

\paragraph{Quotient representatives.}
The $\dimQtenQ$ quotient representatives are indexed by the free columns of the
rank-$\czMinorSize$ certificate of appendix~\ref{app:closure}, and are listed in
table~\ref{tab:quotient}. They are one choice among many: proposition
\ref{prop:cardinality} bounds the cardinality of any closing set from below, but
says nothing about uniqueness, and we make no canonicality claim for this choice.
